% ---------------------------------------------------------------------------
% Appendix: equilibrium after a uniform loss (DERIVED, ideal gas)
% ---------------------------------------------------------------------------
\appendix
\section{Equilibrium of a closed gas after a uniform loss}
\label{app:uniform}
A loss that acts with the same rate on every atom removes a fraction $\delta$
of the atoms without changing the energy per atom or the condensed fraction.
For an ideal Bose gas with a saturated thermal cloud at temperature $T$ in a
harmonic trap, the thermal number and energy are
$N_{\rm th}=\zeta(3)(k_BT/\hbar\bar\omega)^3$ and
$E=3\zeta(4)(k_BT)^4/(\hbar\bar\omega)^3$; the condensate carries a negligible
energy in this approximation. After the loss the gas has $(1-\delta)N$ atoms
and energy $(1-\delta)E$, and if it re-equilibrates at fixed number and energy
its temperature and condensed fraction become
\begin{equation}
T'=(1-\delta)^{1/4}T,\qquad F'=1-(1-\delta)^{-1/4}(1-F).
\label{eq:uniform}
\end{equation}
Here, $F$ is the condensed fraction before the loss. For $\delta=0.35$ and
$F=0.65$, the values of the no-dissipation data at $t=2$~s and of the
reduction of $N$ by the K atoms, Eq.~(\ref{eq:uniform}) gives $F'=0.61$ and
$T'=0.90\,T$, whereas the equilibrium relation at the \emph{unchanged}
temperature gives $F=1-(1-\delta)^{-1}(1-F)=0.46$. Hence a uniform loss
applied to an equilibrium gas leaves it above the constant-temperature
equilibrium curve by $0.15$ in $F$, a difference that disappears when the
comparison is made at fixed energy. The same argument shows that the
re-equilibration proceeds by transferring atoms from the condensate to the
thermal cloud \emph{and} cooling the thermal cloud: the two are not
independent, because each atom leaving the condensate at energy $\mu_c$ has to
be raised to the mean thermal energy $\simeq3k_BT$.
